%=====================================================================
%  Chapter 3 — Hardware and Mechatronic Design
%=====================================================================
\chapter{Hardware and Mechatronic Design}
\label{chap:hardware}

\section{Introduction}
Designing a bipedal robot is a multidisciplinary task that combines
principles from mechanical engineering, electronics, and computer
science. This chapter presents the physical foundation of the
\textbf{Abbes} robot: the mechanical design derived from human
proportions, the selection of actuators and sensors, the power supply,
and the electronic architecture built around an ESP32 microcontroller.
Together these subsystems form the platform on which the firmware and
control software of the following chapters run.

\section{Mechanical Design}
The mechanical design is inspired by the anatomical structure of the
human body. As the review of existing humanoids in
Chapter~\ref{chap:presentation} showed, the construction of humanoid
robots is rooted in mimicking human morphology.

\subsection{Anthropometric Model}
The anthropometric model expresses the robot's dimensions as a function
of human-body proportions, so that its segments and movements stay as
close to human anatomy as possible. Anthropometry --- from the Greek
\emph{anthropos} (human) and \emph{metron} (measure) --- is the study
and measurement of the physical dimensions of the human body. This work
relies on \emph{static} anthropometric measurements, whose implementation
is simpler, using the normalised segmental proportions established by
Winter. Figure~\ref{fig:anthropometric} gives the length of each limb
relative to the total height $H$ of the individual.

\begin{figure}[ht]
  \centering
  \includegraphics[width=0.6\textwidth]{fig_anthropometric.png}
  \caption{Anthropometric model: segment lengths as a proportion of
           total height $H$ (Winter's model).}
  \label{fig:anthropometric}
\end{figure}

\subsection{Mechanical Articulated Structure}
The articulated structure is the set of interconnected joints and links
that lets the robot reproduce human-like movements. The number of
degrees of freedom (DOF) must be large enough for agility and balance,
yet kept minimal to simplify the model. For a leg, two DOF (at the hip
and ankle) enable pendulum-like motion in the sagittal plane, and two
more provide lateral motion in the frontal plane; a fifth DOF at the knee
is required to clear the foot during the swing phase. This leads to five
DOF per leg. Each arm uses three DOF, and the head two, as detailed in
the next section. Figure~\ref{fig:morphology} shows the resulting joint
arrangement.

\begin{figure}[ht]
  \centering
  \includegraphics[width=0.7\textwidth]{fig_morphology.png}
  \caption{Robot morphology: joint arrangement and reference frame
           (P: pitch, R: roll, Y: yaw).}
  \label{fig:morphology}
\end{figure}

\subsection{Final Design: Joint Limits and Resting Positions}
The final design is an 18-DOF humanoid with the following distribution:
\begin{itemize}[itemsep=2pt,topsep=4pt]
  \item \textbf{Each leg --- 5 DOF:} 2 at the hip (pitch and roll),
        1 at the knee (pitch), and 2 at the ankle (pitch and roll).
  \item \textbf{Each arm --- 3 DOF:} 2 at the shoulder (pitch and roll)
        and 1 at the elbow (roll).
  \item \textbf{Head --- 2 DOF:} head yaw and camera pitch.
\end{itemize}

Table~\ref{tab:joints} summarises the complete joint structure. The
angle range and resting position of each joint are expressed in degrees
of servo travel; these exact values are the single source of truth used
by the firmware (the limits are enforced before any command reaches a
servo, see Chapter~\ref{chap:firmware}). The joints are listed in the
firmware's index order, which groups the right-side joints and the head
on the first PWM board and the left-side joints on the second.

\begin{table}[ht]
  \centering
  \caption{Robot joint structure: axes, angular limits and resting
           positions (servo degrees).}
  \label{tab:joints}
  \small
  \setlength{\tabcolsep}{6pt}
  \renewcommand{\arraystretch}{1.12}
  \begin{tabular}{@{}clccccc@{}}
    \toprule
    \textbf{\#} & \textbf{Joint name} & \textbf{Axis} &
    \textbf{Min ($^\circ$)} & \textbf{Max ($^\circ$)} &
    \textbf{Rest ($^\circ$)} \\
    \midrule
    1  & r\_shoulder\_pitch & pitch & 0  & 180 & 90  \\
    2  & r\_shoulder\_roll  & roll  & 0  & 180 & 90  \\
    3  & r\_elbow\_roll     & roll  & 0  & 180 & 90  \\
    4  & head\_yaw          & yaw   & 0  & 180 & 90  \\
    5  & camera\_pitch      & pitch & 0  & 180 & 90  \\
    6  & r\_hip\_roll       & roll  & 0  & 170 & 50  \\
    7  & r\_hip\_pitch      & pitch & 0  & 170 & 80  \\
    8  & r\_knee\_pitch     & pitch & 0  & 170 & 100 \\
    9  & r\_ankle\_pitch    & pitch & 10 & 180 & 100 \\
    10 & r\_ankle\_roll     & roll  & 20 & 150 & 90  \\
    11 & l\_shoulder\_pitch & pitch & 0  & 180 & 90  \\
    12 & l\_shoulder\_roll  & roll  & 0  & 180 & 90  \\
    13 & l\_elbow\_roll     & roll  & 0  & 180 & 90  \\
    14 & l\_hip\_roll       & roll  & 0  & 170 & 50  \\
    15 & l\_hip\_pitch      & pitch & 0  & 170 & 80  \\
    16 & l\_knee\_pitch     & pitch & 0  & 170 & 70  \\
    17 & l\_ankle\_pitch    & pitch & 10 & 180 & 90  \\
    18 & l\_ankle\_roll     & roll  & 20 & 150 & 90  \\
    \bottomrule
  \end{tabular}
\end{table}

The mechanical parts are designed in CAD (Figure~\ref{fig:cad}) and
manufactured by 3D printing; the assembled robot is shown in
Figure~\ref{fig:realrobot}.

\begin{figure}[ht]
  \centering
  \begin{subfigure}{0.40\textwidth}
    \centering
    \includegraphics[height=6.5cm]{fig_cad_design.png}
    \caption{CAD model}
    \label{fig:cad}
  \end{subfigure}
  \hspace{0.8cm}
  \begin{subfigure}{0.40\textwidth}
    \centering
    \includegraphics[height=6.5cm]{fig_real_robot.png}
    \caption{Assembled robot}
    \label{fig:realrobot}
  \end{subfigure}
  \caption{Mechanical design of Abbes: (a) CAD model and (b) the
           3D-printed, assembled robot.}
  \label{fig:mechanical}
\end{figure}

\section{Hardware Components}

\subsection{Actuators}
Selecting actuators involves several criteria: high torque and precision
for accurate motion, efficient current consumption, a size and weight
compatible with the structure, durability, and cost. Table~\ref{tab:act}
compares the candidates considered. The \textbf{MG996R} servo was
selected: it offers a sufficient torque-to-cost ratio for a child-sized,
3D-printed platform, and it is directly compatible with standard PWM
control. The robot uses eighteen such servos, one per joint.

\begin{table}[ht]
  \centering
  \caption{Actuator specifications.}
  \label{tab:act}
  \small
  \renewcommand{\arraystretch}{1.2}
  \begin{tabularx}{\textwidth}{@{}lccccc@{}}
    \toprule
    \textbf{Actuator} & \textbf{Model} & \textbf{Price (\euro)} &
    \textbf{Torque (kg$\cdot$cm)} & \textbf{Position sensor} &
    \textbf{Current (A)} \\
    \midrule
    Servo      & MG996R  & 20    & 10        & --          & 1   \\
    Dynamixel  & MX-28   & 312   & 31.6      & Contactless & 1.7 \\
    Servo      & HS-645MG& 38.40 & 107 / 133 & --          & 2.5 \\
    \bottomrule
  \end{tabularx}
\end{table}

\subsection{Sensors}
The robot's primary sensor is the \textbf{MPU6050}, which combines a
three-axis gyroscope and a three-axis accelerometer in a single
low-cost, I\textsuperscript{2}C-compatible package
(Table~\ref{tab:sensor}). It provides the angular-velocity and
acceleration data needed to estimate the robot's orientation and to
maintain balance during motion. Its digital interface simplifies
integration with the microcontroller.

\begin{table}[ht]
  \centering
  \caption{Sensor description.}
  \label{tab:sensor}
  \small
  \renewcommand{\arraystretch}{1.2}
  \begin{tabularx}{0.7\textwidth}{@{}lX@{}}
    \toprule
    \textbf{Sensor} & \textbf{Description} \\
    \midrule
    MPU6050 & Combines a three-axis gyroscope and a three-axis
              accelerometer (raw I\textsuperscript{2}C). \\
    \bottomrule
  \end{tabularx}
\end{table}

\subsection{Power Supply}
The power supply must deliver the voltage and current required by all
components within budget constraints. The robot is \emph{currently
powered from an off-board source}; a battery-powered version is planned
as future work. Table~\ref{tab:battery} lists the battery options
selected for that version: a high-capacity LiPo pack is preferred for its
current rating, at a higher cost than a basic lithium cell.

\begin{table}[ht]
  \centering
  \caption{Battery specifications.}
  \label{tab:battery}
  \small
  \renewcommand{\arraystretch}{1.2}
  \begin{tabular}{@{}lcc@{}}
    \toprule
    \textbf{Battery} & \textbf{Current rating (mAh)} & \textbf{Cost} \\
    \midrule
    LiPo    & 5000 & High \\
    Lithium & 1000 & Low  \\
    \bottomrule
  \end{tabular}
\end{table}

\subsection{DC/DC Regulator}
A DC/DC converter adapts the battery voltage to the levels required by
the servos and the electronics. Table~\ref{tab:regulator} compares the
regulators considered, trading current rating against cost.

\begin{table}[ht]
  \centering
  \caption{Regulator specifications.}
  \label{tab:regulator}
  \small
  \renewcommand{\arraystretch}{1.2}
  \begin{tabular}{@{}lcc@{}}
    \toprule
    \textbf{Regulator} & \textbf{Current rating (A)} & \textbf{Cost (\euro)} \\
    \midrule
    XL6009 & 3.0 & 14 \\
    XL4016 & 9.0 & 30 \\
    LM2596 & 1.5 & 8  \\
    \bottomrule
  \end{tabular}
\end{table}

\section{Electronic Architecture}
\label{sec:hw-arch}
The robot's electronics are organised into a \textbf{control section} and
a \textbf{power section}.

The control section is built around an \textbf{ESP32 microcontroller},
which runs the real-time control loop directly on the robot and joins the
ROS\,2 graph as a \emph{micro-ROS} node over WiFi/UDP. The ESP32 drives
the eighteen servos through \textbf{two PCA9685} 16-channel PWM driver
boards on the I\textsuperscript{2}C bus (at addresses 0x40 and 0x41): the
first board carries the right-side and head joints, the second the
left-side joints. The same I\textsuperscript{2}C bus connects the
\textbf{MPU6050} IMU. On the host (laptop) side, a micro-ROS \emph{agent}
bridges the ESP32 into ROS\,2, so that the inertial data published by the
node and the joint commands it receives are available as ordinary ROS\,2
topics. This architecture replaces the more common pairing of a
single-board computer with a separate servo controller by a single,
low-cost microcontroller that communicates wirelessly. The power section
supplies the servos and electronics through the battery and DC/DC
regulator described above. Figure~\ref{fig:hwarch} summarises the
architecture.

\begin{figure}[ht]
  \centering
  \begin{tikzpicture}[
      font=\small,
      box/.style={draw, rounded corners=2pt, align=center,
                  minimum height=9mm, minimum width=22mm, thick},
      host/.style={box, fill=blue!8},
      mcu/.style={box, fill=orange!12},
      drv/.style={box, fill=green!10},
      sen/.style={box, fill=purple!8},
      pwr/.style={box, fill=red!8},
      >=Latex
    ]
    \node[host] (host) {Host PC\\\scriptsize ROS\,2 + micro-ROS agent};
    \node[mcu, right=22mm of host] (esp) {ESP32\\\scriptsize micro-ROS node};
    \node[drv, right=20mm of esp, yshift=12mm] (pca0) {PCA9685 \#0\\\scriptsize 0x40};
    \node[drv, right=20mm of esp, yshift=-12mm] (pca1) {PCA9685 \#1\\\scriptsize 0x41};
    \node[sen, below=10mm of esp] (imu) {MPU6050\\\scriptsize IMU};
    \node[box, right=14mm of pca0, fill=gray!10] (s0) {10 servos\\\scriptsize joints 1--10};
    \node[box, right=14mm of pca1, fill=gray!10] (s1) {8 servos\\\scriptsize joints 11--18};
    \node[pwr, below=10mm of imu, xshift=20mm] (pwr) {Battery $\rightarrow$ DC/DC regulator};

    \draw[<->] (host) -- node[above,font=\scriptsize]{WiFi/UDP} (esp);
    \draw[->] (esp) -- node[above,sloped,font=\scriptsize]{I\textsuperscript{2}C} (pca0);
    \draw[->] (esp) -- node[below,sloped,font=\scriptsize]{I\textsuperscript{2}C} (pca1);
    \draw[<->] (esp) -- node[right,font=\scriptsize]{I\textsuperscript{2}C} (imu);
    \draw[->] (pca0) -- (s0);
    \draw[->] (pca1) -- (s1);
    \draw[-, dashed, red!60] (pwr) -- (esp);
    \draw[-, dashed, red!60] (pwr) -| (s0);
    \draw[-, dashed, red!60] (pwr) -| (s1);
  \end{tikzpicture}
  \caption{Electronic architecture of Abbes: the ESP32 micro-ROS node
           drives 18 servos through two PCA9685 boards and reads the
           MPU6050, communicating with the ROS\,2 host over WiFi. Dashed
           red lines indicate power distribution.}
  \label{fig:hwarch}
\end{figure}

\section{Conclusion}
This chapter presented the hardware foundation of Abbes. The mechanical
structure was derived from anthropometric data to obtain a human-like,
18-DOF morphology, and its joint limits were defined
(Table~\ref{tab:joints}). The components were selected to balance
performance, cost, and efficiency: MG996R servos, an MPU6050 IMU, and a
battery/regulator pair for the powered version. The electronic
architecture is built around an ESP32 that drives the servos through two
PCA9685 boards and joins ROS\,2 as a micro-ROS node. The next chapter
details the embedded firmware that runs on this ESP32.
